% ============================================================
% Chapter 7: Safety Checklist for Retrofits
% ============================================================
\chapter{Safety Checklist for Retrofits}
\label{chap:safety_checklist}

This chapter consolidates practical artefacts derived from the project's verification activities (\Cref{chap:testing}) and from principles of relevant functional safety standards (\Cref{sec:standards_detail}). It presents a checklist of verification points for this retrofit project, each marked with the state it reached, as an example for verifying similar modernisations of safety-critical electronics.

% ------------------------------------------------------------
\section{Safety Verification Checklist for Retrofits}
\label{sec:safety_checklist_items}

This checklist consolidates key verification activities based on the methodology (\Cref{chap:testing}) and principles drawn from standards such as~\gls{iec61508}~\cite{iec61508_2010}, \gls{en50129}~\cite{en50129_2018}, \gls{en50716}~\cite{en50716_2023}, and~\gls{ieee1012}~\cite{ieee1012_2024}. It serves as a guide for verifying similar safety-critical electronic retrofits, illustrating the types of checks across the lifecycle. The status markers reflect the state reached \textit{within this specific research project}, as documented in this report: \texttt{[DONE]} where the report documents the item as complete, \texttt{[PARTIAL]} where it documents part of it, \texttt{[OPEN]} where it documents nothing that meets the item, and \texttt{[PLANNED]} for future work (\Cref{sec:future_work}); items marked \texttt{[PARTIAL]} or \texttt{[OPEN]} give the reason. The items are numbered within their group and carry the prefix CL (CL-A1, CL-A2, \dots), so that the text can refer to them without confusing them with the defects B1 to B6 of \Cref{tab:bug_taxonomy}. The evidence behind the markers is summarised in \Cref{tab:evidence_state}.

\bigskip

\noindent\textbf{A. Requirements, Design Documentation, and Planning}
\begin{checklist}
	\chk{A1}{PARTIAL} Legacy functional requirements: taken from the legacy manual~\cite{rexroth1992wsp} (diagnosis table, operator functions) and from the prior reverse engineering~\cite{ayoub2014, leyer2023} (\Cref{sec:derivation_of_strategy,sec:system_scenario_testing}); no consolidated requirements specification is part of this report.
	\chk{A2}{PARTIAL} Requirements for new functionalities (enhanced diagnostics via \gls{stm32}): the functions are described and kept apart from the safety functions by design (\Cref{sec:stm32_details}); they are not specified as requirements, and the diagnostic frame differs between the designs of 2023 and 2025 (\Cref{sec:fpga_core_details}).
	\chk{A3}{PARTIAL} Modifications introduced by the retrofit, compared to the original design: described with the architecture (\Cref{chap:architecture}), including two that change the behaviour of the module, the added brown-out detection and the~\gls{nvram} that needs no STORE and RECALL signals, so that the STORE and RECALL inputs of \texttt{X4} no longer act on the fault memory (\Cref{sec:en50716_modifications}); no itemised comparison with the original design and no documented impact analysis (\Cref{sec:en50716_modifications}).
	\chk{A4}{OPEN} \Gls{hra} of legacy functions and new components: no~\gls{fmea}, \gls{pha} or other~\gls{hra} is documented, and no hazard is identified; \Cref{sec:results_discussion_risk} discusses technical risks of the project qualitatively, and \Cref{sec:vehicle_outputs} names an open question for such an analysis: whether a reset in service leads into the automatic check.
	\chk{A5}{PARTIAL} Mapping to applicable standards: the verification layers are mapped to whole standards, with pointers to the parts and annexes the text relies on (\Cref{tab:standards_mapping}); requirements and design changes are not mapped to individual clauses, which needs the full texts of the standards (\Cref{sec:standards_detail,sec:safety_compliance}).
	\chk{A6}{PARTIAL} Verification and validation (V\&V) approach described in~\Cref{chap:testing} (scope, Layer~1 and Layer~2, tools); pass/fail criteria are given only as the expected codes of the~\gls{yaml} test cases (\Cref{sec:app_yaml_test_cases}).
	\chk{A7}{PARTIAL} Traceability (linking requirements, design, code, tests, results): no requirement list exists; the Layer~2 cases are linked to the faults of the diagnosis table by name, the Layer~1 tests only to the functions identified by reverse engineering, and no traceability matrix exists (\Cref{sec:verification_artefacts}).
	\chk{A8}{PARTIAL} Configuration management: the follow-on work kept its sources in Git~\cite[p.~7]{savchuk2025}; the Layer~2 test definitions and logs are not archived (\Cref{sec:verification_artefacts}).
\end{checklist}

\medskip
\noindent\textbf{B. Hardware Component and Interface Verification (Layer~1)} (\Cref{sec:fpga_interface_testing})
\begin{checklist}
	\chk{B1}{PARTIAL} \gls{vhdl} logic modules exercised in the simulations of 2023 using dedicated testbenches; the 17 archived testbench files of 2023 contain no assertions (\Cref{sec:fpga_interface_testing}).
	\chk{B2}{PARTIAL} Simulation results reviewed and documented: for the design of 2023, checked by manual waveform analysis~\cite[Sec.~5.1]{leyer2023} (\Cref{sec:fpga_interface_testing,subsec:fpga_interface_results}).
	\chk{B3}{PARTIAL} Memory interfaces (Flash~\gls{rom}, \gls{nvram}) exercised in simulation and with test designs on the prototype; the simulation could not detect defect~B1, because its model of the \gls{nvram} ignores the byte-low enable, and no simulation after the correction is archived (\Cref{subsec:fpga_interface_results}).
	\chk{B4}{PARTIAL} Critical peripheral logic (display driving, button handling, frequency sweep, \gls{spi}~master, watchdog interaction) exercised in the simulations of 2023 (\Cref{sec:fpga_interface_testing}, \Cref{tab:evidence_state}).
	\chk{B5}{PARTIAL} \gls{fpga} identification and programming verified via~\gls{jtag}, and pin connectivity checked with a toggle test, on the Rev~0.0 prototype in 2023~\cite{leyer2023}; the solder defects (B3) surfaced later, during commissioning of the rebuilt unit, and no \gls{jtag} check or toggle test after the re-soldering of the \gls{fpga} pins is documented (\Cref{subsec:fpga_interface_results}).
	\chk{B6}{PARTIAL} Critical hardware timing parameters (watchdog behaviour, frequency sweep) measured on the inherited prototype~\cite{leyer2023} (\Cref{fig:watchdog_reaction,fig:sweep_comparison}); on the reworked assembly, no measurement of either is documented, nor any observation of the watchdog after the re-soldering of the \gls{fpga} pins (\Cref{subsec:fpga_interface_results}).
\end{checklist}

\medskip
\noindent\textbf{C. Software/Firmware Verification (as applicable)}
\begin{checklist}
	\chk{C1}{PARTIAL} \gls{stm32} firmware modules (SPI handler, RTC management, SD card logging, USB interface): fix applied to the firmware of 2023 for the SPI reception and FatFs defects; diagnostic path not verified separately (\Cref{sec:diagnostic_resolution}).
	\chk{C2}{OPEN} Verification of adherence to defined coding standards: no coding standard is defined in this report for the~\gls{vhdl} or the firmware.
	\chk{C3}{PARTIAL} Legacy firmware kept unchanged by design (\Cref{sec:fpga_core_details}); the Flash~\gls{rom} was programmed and verified on the Rev~0.0 prototype in 2023~\cite[Sec.~4.3]{leyer2023}, but no checksum and no comparison with the contents of the original~\gls{eprom} are reported; its execution within the~\gls{fpga} emulation core observed in the follow-on fixture runs (\Cref{subsec:equivalence_discussion}).
\end{checklist}

\medskip
\noindent\textbf{D. Integration and System-Level Testing (Layer~2)} (\Cref{sec:system_scenario_testing})
\begin{checklist}
	\chk{D1}{PARTIAL} Test scenarios for the documented legacy fault handling~\cite{rexroth1992wsp} defined in the~\gls{yaml} format of the fixture~\cite{dolata2023} (example in~\Cref{app:test_cases}); the check of the self-test was still outstanding in the fixture of 2023~\cite[Sec.~6.2]{dolata2023}, and the definitions used in the follow-on runs are not archived (\Cref{sec:verification_artefacts}).
	\chk{D2}{DONE} Tests executed on integrated hardware using the dedicated test fixture~\cite{dolata2023}, corrected by the follow-on team~\cite[pp.~10--15]{savchuk2025}.
	\chk{D3}{PARTIAL} Fault memory cycle exercised on the fixture after the correction of defect~B1 (detection, display, STW read-out, clearing); storage across power cycles not documented (\Cref{subsec:system_scenario_results}).
	\chk{D4}{OPEN} Legacy self-test sequences: the scenario is defined in the method against the documentation~\cite{rexroth1992wsp} (\Cref{sec:system_scenario_testing}), but its execution on the reworked assembly is not documented; only the start of the sequence was shown, in Leyer's simulation of 2023 (\Cref{subsec:system_scenario_results}).
	\chk{D5}{PARTIAL} User interface behaviour (button responses, display outputs): STW and Lösch exercised in the fixture runs; Test and Tür/V not documented (\Cref{subsec:system_scenario_results}).
	\chk{D6}{PARTIAL} Enhanced diagnostic logging functionality (\gls{stm32} subsystem receiving data via isolated~\gls{spi}): fix applied for defect B4 and for the receiver part of B5, both in the firmware of 2023; diagnostic path not verified separately; for the transmitter part of B5 and the configuration of the rebuilt unit see \Cref{sec:diagnostic_resolution}.
	\chk{D7}{PARTIAL} System stability: the rebuilt unit was observed operating as a whole in the test setup; stability under power cycling not documented (\Cref{subsec:system_scenario_results}).
	\chk{D8}{OPEN} Test results logged, reviewed and archived: only the log format of the fixture application is documented (\Cref{sec:app_system_log_excerpts}).
\end{checklist}

\medskip
\noindent\textbf{E. Environmental and EMC Considerations}
\begin{checklist}
	% a smaller emergency stretch than the document's sets this item without loose lines
	\chk{E1}{PARTIAL} {\setlength{\emergencystretch}{1em}Design principles with a view to \gls{en50155} and \gls{en61373} (component selection, layout, isolation) are described in~\Cref{sec:design_principles}, which cites these standards by title only; layout and isolation are documented (\Cref{sec:pcb_design_impl}), but no component is documented as chosen for, or checked against, these standards, and the~\gls{bom} (\Cref{app:bom_drc}) has no column for temperature or vibration ratings. For the~\gls{fpga} it gives the part number \texttt{A3P1000-PQG208M}; in the ordering information of the manufacturer's datasheet for the military ProASIC3/EL devices~\cite[Fig.~67]{microsemi_ds0106}, which is archived with Leyer's documentation, the suffix~M stands for the military temperature range, \SIrange{-55}{125}{\celsius} junction temperature. Leyer calls these~\glspl{fpga} \enquote{military graded} and writes that, for lack of new ones, they were desoldered from broken boards~\cite[Sec.~6.2]{leyer2023}, so the Base Board carries a used part.\par}
	\chk{E2}{PARTIAL} \gls{emc} mitigation techniques implemented and documented: the design measures that bear on~\gls{emc} (ground planes, decoupling, galvanic isolation of the~\gls{spi} link) are described (\Cref{sec:pcb_design_impl}; stack-up in~\Cref{sec:app_drc}), but their effect has not been tested: no~\gls{emc} measurement or test of the rebuilt unit is documented, and the galvanic isolation has not been tested separately (\Cref{sec:results_discussion_risk}).
	\chk{E3}{PLANNED} Qualification tests for the environment (\gls{en50155}, \gls{en61373}) and for~\gls{emc} (\gls{en50121}~series), to be performed by an accredited laboratory for a formal certification (\Cref{sec:future_work}).
\end{checklist}

\medskip
\noindent\textbf{F. Manufacturing and Documentation}
\begin{checklist}
	\chk{F1}{PARTIAL} \gls{pcb} manufacturing outputs inspected (visual checks, basic electrical continuity/short checks where feasible): the follow-on project reports that the diagnostics board showed no short circuit~\cite[p.~22]{savchuk2025}; no other inspection record is documented.
	\chk{F2}{PARTIAL} Assembly guided by the Rev~1.0 schematic and~\gls{bom} (\Cref{sec:pcb_design_impl}); apart from the inspection of the populated boards for the~\gls{bom} (\Cref{app:bom_drc}), no inspection of the rebuilt unit before its commissioning is documented, and the solder defects on its~\gls{fpga} (B3) surfaced only during commissioning (\Cref{sec:integration_defects}). The inspection that found two cold solder joints was made on the Rev~0.0 prototype in 2023~\cite[Sec.~4.3]{leyer2023} (\Cref{sec:fpga_interface_testing}).
	\chk{F3}{PARTIAL} Final design documentation (\gls{bom} in~\Cref{app:bom_drc}, schematic and layouts in~\Cref{app:schematics}) for the reworked assembly (Rev~1.0 Upper, LED and \gls{stm32} boards on the Rev~0.0 Base Board): the schematic comes from the Rev~1.0 file while the fitted Base Board is Rev~0.0, and the notes to~\Cref{tab:bom} record where the sources disagree, so it is not an \enquote{as-built} record in every detail.
	\chk{F4}{PARTIAL} Key verification artefacts archived: design files and simulation sources; the Layer~2 test definitions and logs of the follow-on runs are not archived (\Cref{app:test_logs,app:test_cases}).
	\chk{F5}{PARTIAL} Traceability information: the Layer~2 cases are linked to the legacy faults by name; there is no requirement list and no traceability matrix (\Cref{sec:verification_artefacts}).
	\chk{F6}{PARTIAL} Development and verification tools named (\Cref{sec:verification_artefacts}); their classification and a justification of their suitability were beyond the project's scope.
\end{checklist}